\section*{Wykaz skrótów}
\addcontentsline{toc}{section}{Wykaz skrótów}

\begin{tabular}{ll}
\textbf{AI} & Artificial Intelligence (sztuczna inteligencja) \\
\textbf{API} & Application Programming Interface (interfejs programistyczny) \\
\textbf{AOI} & Area of Interest (obszar zainteresowania) \\
\textbf{CEMS} & Copernicus Emergency Management Service \\
\textbf{CDS} & Climate Data Store \\
\textbf{GeoJSON} & Geographic JavaScript Object Notation \\
\textbf{GeoTIFF} & Geographic Tagged Image File Format \\
\textbf{GIS} & Geographic Information System (system informacji geograficznej) \\
\textbf{GloFAS} & Global Flood Awareness System \\
\textbf{HTTP} & Hypertext Transfer Protocol \\
\textbf{LLM} & Large Language Model (duży model językowy) \\
\textbf{OSM} & OpenStreetMap \\
\textbf{POI} & Point of Interest (punkt zainteresowania) \\
\textbf{REST} & Representational State Transfer \\
\textbf{RP} & Return Period (okres powrotny) \\
\textbf{UUID} & Universally Unique Identifier \\
\end{tabular}

\newpage
\section*{Wykaz definicji}
\addcontentsline{toc}{section}{Wykaz definicji}

\begin{description}
    \item[Agent AI] Program komputerowy wykorzystujący duży model językowy (LLM) do autonomicznego wykonywania zadań poprzez wybór i~sekwencjonowanie dostępnych narzędzi.

    \item[AgentOS] Wzorzec architektoniczny z~frameworka Agno, który automatycznie generuje aplikację FastAPI z~endpointami dla agenta AI.

    \item[Analiza przestrzenna] Metoda analizy danych geograficznych polegająca na badaniu relacji przestrzennych między obiektami (np. przecięcia, odległości, zawieranie).

    \item[Bounding box] Prostokątny obszar definiowany przez współrzędne narożników (minX, minY, maxX, maxY), używany do filtrowania danych przestrzennych.

    \item[Geokodowanie] Proces konwersji adresów lub nazw miejscowości na współrzędne geograficzne.

    \item[GloFAS] Globalny system wczesnego ostrzegania przed powodziami prowadzony przez Copernicus Emergency Management Service.

    \item[Okres powrotny] Statystyczny czas między zdarzeniami o~określonej skali (np. RP100 oznacza powódź występującą średnio raz na 100~lat).

    \item[Overture Maps] Otwarty zbiór danych geograficznych tworzony przez konsorcjum firm technologicznych, zawierający budynki, POI i~infrastrukturę.

    \item[Spatial Join] Operacja łączenia zbiorów danych geograficznych na podstawie relacji przestrzennych (np. przecięcia geometrii).

    \item[Tool (narzędzie)] Funkcja Python z~dekoratorem \texttt{@tool}, którą agent AI może wywołać w~celu wykonania określonego zadania.
\end{description}

\newpage
\section*{Wykaz załączników}
\addcontentsline{toc}{section}{Wykaz załączników}

\begin{description}
    \item[Załącznik A] Kod źródłowy systemu RiskScoop AI
    \begin{itemize}
        \item A.1 -- Plik główny \texttt{my\_os.py} (konfiguracja AgentOS)
        \item A.2 -- Narzędzia agenta (\texttt{tools/*.py})
        \item A.3 -- Serwisy biznesowe (\texttt{services/*.py})
        \item A.4 -- Frontend (\texttt{test\_frontend/index.html})
    \end{itemize}

    \item[Załącznik B] Screenshoty interfejsu użytkownika
    \begin{itemize}
        \item B.1 -- Analiza szpitali w~Genewie (rys.~\ref{fig:example_geneva})
        \item B.2 -- Analiza szkół w~Warszawie (rys.~\ref{fig:example_warsaw})
        \item B.3 -- Analiza budynków w~Krakowie (rys.~\ref{fig:example_krakow})
    \end{itemize}
\end{description}

